\section{预备知识与问题定义}

\subsection{基础概念与符号说明}

本章用于补充读者理解后续方案所需的基础知识，例如密码学原语、访问控制模型、流量分析概念、攻击检测指标或协议结构等。建议仅保留与论文核心贡献直接相关的内容，避免将大段教材式介绍堆叠进正文。

\subsection{问题描述与研究目标}

本节应明确论文所研究的问题边界、输入输出、适用场景和评价目标。若论文包含形式化描述，可在此给出模型、符号和约束条件，并说明研究对象与已有工作相比的差异。

\begin{equation}
  \mathcal{L} = \sum_{i=1}^{n} \omega_i \cdot \ell_i
\end{equation}

上式仅作为公式排版示例。正式写作时，建议在公式前后分别交代变量含义、适用条件和推导作用，并确保编号连续、引用一致。

\subsection{需求分析}

若论文包含系统设计或工程实现，可在本节整理功能需求、性能需求、安全需求和可用性要求等内容。对于关键约束条件，可使用分项表述，但仍应保持学术论文的叙述风格，而非产品文档式表达。
